% =====================================================================
%  modelos_n_gramas.tex
%  ========================
%
%
%  Descripción:
%  ------------
%  Material de apoyo sobre modelos de lenguaje de n-gramas: la regla de
%  la cadena, el supuesto de Markov, la estimación por máxima
%  verosimilitud, el suavizado y la perplejidad. Va aparte de la
%  libreta porque no es toma de clase sino desarrollo de fondo para
%  repasar antes del examen.
%
%
%  Considerations:
%  ------------
%  - El texto lo generó Claude, no salió de la clase. Lo que quedó sin
%    verificar contra el material del curso va marcado como tal en el
%    propio documento, en los bloques de Observación
%  - Carga el preámbulo compartido, así que usa los mismos entornos de
%    teorema y las mismas cajas de color que las libretas. Las cajas de
%    intuición son la caja idea con su título propio
%  - Los macros de notación (\Count, \Vocab, \bos, \eos, \unk) se
%    declaran después del \input porque solo los usa este archivo
%  - No lleva portadilla: el documento no es de autoría propia, así que
%    no se firma como las libretas
%
%
%  Metadata:
%  ----------
%   * Author : claude
%   * Version: 1.0
%
%
%  History:
%  ------------
%  Author      Date            Description
%  claude       22/08/2026      Creation
%
%
% =====================================================================
\documentclass[11pt,letterpaper]{article}

% =====================================================================
%  DATOS DE LA MATERIA
% =====================================================================
\newcommand{\materia}{Fundamentos para el Análisis Automático de Lenguaje}
\newcommand{\profesor}{Dra. Delia Irazú Hernández Farías \\
                       Dr. Luis Villaseñor Pineda \\
                       Dr. Manuel Montes y Gómez \\
                       Dr. Aurelio López López}
\newcommand{\periodo}{Otoño 2026}
\newcommand{\alumno}{Bryan Ezequiel Violante Arriaga}

% =====================================================================
%  Preambulo.tex
%  ========================
%
%
%  Descripción:
%  ------------
%  Preámbulo compartido de los apuntes de clase: paquetes, estilo de
%  página, entornos de teorema, las cuatro cajas de color, el estilo de
%  listings y el pseudocódigo en español. Se carga con \input desde
%  cada documento en lugar de copiarlo, para que un arreglo aquí llegue
%  a todas las libretas de una vez.
%
%
%  Considerations:
%  ------------
%  - Se carga con \input y ruta relativa, no con \usepackage, porque
%    \usepackage solo busca en la carpeta del documento y en el árbol
%    de TeX. Un .sty instalado en TEXMFHOME evitaría la ruta, pero vive
%    fuera del repo y una clona nueva no compilaría
%  - Los cuatro datos de la materia se declaran en el documento ANTES
%    del \input. hyperref expande \materia al cargarse, así que si no
%    existe todavía la compilación truena
%  - Cuando la materia la imparte más de un profesor, los nombres van
%    en un solo \profesor separados por \\. La portadilla los imprime
%    dentro de un center, que es donde \\ sí corta renglón. Repetir
%    \newcommand{\profesor} no es opción: truena con "already defined"
%  - algorithmicx no tiene opción de idioma y babel no lo traduce, así
%    que las palabras clave del pseudocódigo se renombran a mano
%  - Los nombres de los entornos van sin acento (\begin{definicion})
%    aunque el rótulo impreso sí lo lleve
%  - babel se carga con es-noshorthands porque en español el " queda
%    activo y se come el texto que le sigue, además de romper listings
%
%
%  Metadata:
%  ----------
%   * Author : zxxz6 (Bryan Violante Arriaga)
%   * Version: 1.7.0
%
%
%  History:
%  ------------
%  Author      Date            Description
%  zxxz6       24/09/2026      El nombre del alumno en el encabezado
%  zxxz6       24/09/2026      Encabezado de una hoja de un ejercicio
%  zxxz6       23/09/2026      Macros de los documentos de soluciones
%  zxxz6       30/08/2026      Macros de lim, limsup y liminf sobre n
%  zxxz6       25/08/2026      Macros de Omega, Theta, o y omega chicas
%  zxxz6       20/08/2026      Cargue pgfplots para graficas de datos
%  zxxz6       20/08/2026      Cargue tikz para diagramas de conjuntos
%  zxxz6       19/08/2026      Documente \profesor con varios nombres
%  zxxz6       19/08/2026      Agregue la caja Idea Random
%  zxxz6       18/08/2026      Creation
%
%
% =====================================================================
%  USO
%  ---
%  \documentclass[11pt,letterpaper]{article}
%  \newcommand{\materia}{...}      % los cuatro datos van arriba
%  \newcommand{\profesor}{...}
%  \newcommand{\periodo}{...}
%  \newcommand{\alumno}{...}
%  % =====================================================================
%  Preambulo.tex
%  ========================
%
%
%  Descripción:
%  ------------
%  Preámbulo compartido de los apuntes de clase: paquetes, estilo de
%  página, entornos de teorema, las cuatro cajas de color, el estilo de
%  listings y el pseudocódigo en español. Se carga con \input desde
%  cada documento en lugar de copiarlo, para que un arreglo aquí llegue
%  a todas las libretas de una vez.
%
%
%  Considerations:
%  ------------
%  - Se carga con \input y ruta relativa, no con \usepackage, porque
%    \usepackage solo busca en la carpeta del documento y en el árbol
%    de TeX. Un .sty instalado en TEXMFHOME evitaría la ruta, pero vive
%    fuera del repo y una clona nueva no compilaría
%  - Los cuatro datos de la materia se declaran en el documento ANTES
%    del \input. hyperref expande \materia al cargarse, así que si no
%    existe todavía la compilación truena
%  - Cuando la materia la imparte más de un profesor, los nombres van
%    en un solo \profesor separados por \\. La portadilla los imprime
%    dentro de un center, que es donde \\ sí corta renglón. Repetir
%    \newcommand{\profesor} no es opción: truena con "already defined"
%  - algorithmicx no tiene opción de idioma y babel no lo traduce, así
%    que las palabras clave del pseudocódigo se renombran a mano
%  - Los nombres de los entornos van sin acento (\begin{definicion})
%    aunque el rótulo impreso sí lo lleve
%  - babel se carga con es-noshorthands porque en español el " queda
%    activo y se come el texto que le sigue, además de romper listings
%
%
%  Metadata:
%  ----------
%   * Author : zxxz6 (Bryan Violante Arriaga)
%   * Version: 1.7.0
%
%
%  History:
%  ------------
%  Author      Date            Description
%  zxxz6       24/09/2026      El nombre del alumno en el encabezado
%  zxxz6       24/09/2026      Encabezado de una hoja de un ejercicio
%  zxxz6       23/09/2026      Macros de los documentos de soluciones
%  zxxz6       30/08/2026      Macros de lim, limsup y liminf sobre n
%  zxxz6       25/08/2026      Macros de Omega, Theta, o y omega chicas
%  zxxz6       20/08/2026      Cargue pgfplots para graficas de datos
%  zxxz6       20/08/2026      Cargue tikz para diagramas de conjuntos
%  zxxz6       19/08/2026      Documente \profesor con varios nombres
%  zxxz6       19/08/2026      Agregue la caja Idea Random
%  zxxz6       18/08/2026      Creation
%
%
% =====================================================================
%  USO
%  ---
%  \documentclass[11pt,letterpaper]{article}
%  \newcommand{\materia}{...}      % los cuatro datos van arriba
%  \newcommand{\profesor}{...}
%  \newcommand{\periodo}{...}
%  \newcommand{\alumno}{...}
%  % =====================================================================
%  Preambulo.tex
%  ========================
%
%
%  Descripción:
%  ------------
%  Preámbulo compartido de los apuntes de clase: paquetes, estilo de
%  página, entornos de teorema, las cuatro cajas de color, el estilo de
%  listings y el pseudocódigo en español. Se carga con \input desde
%  cada documento en lugar de copiarlo, para que un arreglo aquí llegue
%  a todas las libretas de una vez.
%
%
%  Considerations:
%  ------------
%  - Se carga con \input y ruta relativa, no con \usepackage, porque
%    \usepackage solo busca en la carpeta del documento y en el árbol
%    de TeX. Un .sty instalado en TEXMFHOME evitaría la ruta, pero vive
%    fuera del repo y una clona nueva no compilaría
%  - Los cuatro datos de la materia se declaran en el documento ANTES
%    del \input. hyperref expande \materia al cargarse, así que si no
%    existe todavía la compilación truena
%  - Cuando la materia la imparte más de un profesor, los nombres van
%    en un solo \profesor separados por \\. La portadilla los imprime
%    dentro de un center, que es donde \\ sí corta renglón. Repetir
%    \newcommand{\profesor} no es opción: truena con "already defined"
%  - algorithmicx no tiene opción de idioma y babel no lo traduce, así
%    que las palabras clave del pseudocódigo se renombran a mano
%  - Los nombres de los entornos van sin acento (\begin{definicion})
%    aunque el rótulo impreso sí lo lleve
%  - babel se carga con es-noshorthands porque en español el " queda
%    activo y se come el texto que le sigue, además de romper listings
%
%
%  Metadata:
%  ----------
%   * Author : zxxz6 (Bryan Violante Arriaga)
%   * Version: 1.7.0
%
%
%  History:
%  ------------
%  Author      Date            Description
%  zxxz6       24/09/2026      El nombre del alumno en el encabezado
%  zxxz6       24/09/2026      Encabezado de una hoja de un ejercicio
%  zxxz6       23/09/2026      Macros de los documentos de soluciones
%  zxxz6       30/08/2026      Macros de lim, limsup y liminf sobre n
%  zxxz6       25/08/2026      Macros de Omega, Theta, o y omega chicas
%  zxxz6       20/08/2026      Cargue pgfplots para graficas de datos
%  zxxz6       20/08/2026      Cargue tikz para diagramas de conjuntos
%  zxxz6       19/08/2026      Documente \profesor con varios nombres
%  zxxz6       19/08/2026      Agregue la caja Idea Random
%  zxxz6       18/08/2026      Creation
%
%
% =====================================================================
%  USO
%  ---
%  \documentclass[11pt,letterpaper]{article}
%  \newcommand{\materia}{...}      % los cuatro datos van arriba
%  \newcommand{\profesor}{...}
%  \newcommand{\periodo}{...}
%  \newcommand{\alumno}{...}
%  \input{../../templates/Preambulo.tex}
%  \begin{document}
%  \portadilla
%
%  Con dos o mas profesores se declara UN solo \profesor y los nombres
%  se separan con \\, uno por renglon:
%
%  \newcommand{\profesor}{Dra. Fulana de Tal \\
%                         Dr. Mengano de Cual}
% =====================================================================


% ==================== IDIOMA Y CODIFICACIÓN ====================
\usepackage[utf8]{inputenc}
\usepackage[T1]{fontenc}
\usepackage[spanish,es-tabla,es-noshorthands]{babel}
\usepackage{lmodern}
\usepackage{microtype}

% ==================== PÁGINA ====================
\usepackage[letterpaper,margin=2.3cm,headheight=14pt]{geometry}
\usepackage{parskip}
\usepackage{fancyhdr}
\usepackage{enumitem}

% ==================== MATEMÁTICAS ====================
\usepackage{amsmath,amssymb,amsthm,mathtools}

% ==================== FIGURAS, TABLAS, CÓDIGO ====================
\usepackage{graphicx}
\usepackage{booktabs}
\usepackage{array}
% float da el especificador [H], que clava una tabla donde se
% escribio. Una traza que LaTeX manda dos paginas adelante deja
% de leerse junto al paso que la produjo
\usepackage{float}
\usepackage{xcolor}
% tikz ya viene arrastrado por tcolorbox[most], pero se carga
% explicito para no depender de ese detalle; la libreria babel
% protege a tikz de los caracteres activos del español
\usepackage{tikz}
\usetikzlibrary{babel}
% pgfplots dibuja graficas de datos (ejes, escalas log) sobre tikz
\usepackage{pgfplots}
\pgfplotsset{compat=1.18}
\usepackage{listings}
\usepackage[ruled,section]{algorithm}
\usepackage{algpseudocode}
\usepackage[most]{tcolorbox}
% soul da \hl, el unico resaltado que parte renglon; \colorbox no
\usepackage{soul}
\usepackage{hyperref}

% =====================================================================
%  ESTILO
% =====================================================================

% --- encabezado y pie de cada página ---
\pagestyle{fancy}
\fancyhf{}
\fancyhead[L]{\small\textsc{\materia}}
\fancyhead[R]{\small\periodo}
\fancyfoot[C]{\small\thepage}
\renewcommand{\headrulewidth}{0.4pt}

% --- enlaces ---
\hypersetup{
  colorlinks=true,
  linkcolor=black,
  urlcolor=blue,
  citecolor=black,
  pdftitle={Apuntes de \materia},
  pdfauthor={\alumno}
}

% --- paleta ---
\definecolor{acento}{HTML}{1F4E79}
\definecolor{fondoIdea}{HTML}{EAF2FB}
\definecolor{fondoOjo}{HTML}{FDF0E6}
\definecolor{fondoPend}{HTML}{F2F0F7}
\definecolor{comentario}{HTML}{5A7D5A}
\definecolor{cadena}{HTML}{A03030}
\definecolor{fondoCodigo}{HTML}{F7F7F7}
\definecolor{fondoResultado}{HTML}{EAF5EC}
\definecolor{marcador}{HTML}{FFF3A3}

% --- entornos de teorema; numeran por sección, o sea por clase ---
\theoremstyle{definition}
\newtheorem{definicion}{Definición}[section]
\newtheorem{ejemplo}[definicion]{Ejemplo}
\newtheorem{ejercicio}[definicion]{Ejercicio}

\theoremstyle{plain}
\newtheorem{teorema}[definicion]{Teorema}
\newtheorem{lema}[definicion]{Lema}
\newtheorem{corolario}[definicion]{Corolario}
\newtheorem{proposicion}[definicion]{Proposición}

\theoremstyle{remark}
\newtheorem*{observacion}{Observación}
\newtheorem*{quick_sum}{Resumen rápido}

% --- cajas de color para lo que se relee antes del examen ---
\newtcolorbox{idea}[1][]{
  colback=fondoIdea, colframe=acento, boxrule=0.6pt,
  arc=2pt, left=6pt, right=6pt, top=5pt, bottom=5pt,
  title={Idea clave}, fonttitle=\bfseries\small, #1
}
\newtcolorbox{ojo}[1][]{
  colback=fondoOjo, colframe=orange!70!black, boxrule=0.6pt,
  arc=2pt, left=6pt, right=6pt, top=5pt, bottom=5pt,
  title={Ojito...}, fonttitle=\bfseries\small, #1
}
\newtcolorbox{pendiente}[1][]{
  colback=fondoPend, colframe=violet!60!black, boxrule=0.6pt,
  arc=2pt, left=6pt, right=6pt, top=5pt, bottom=5pt,
  title={Pendiente por resolver}, fonttitle=\bfseries\small, #1
}
\newtcolorbox{idear}[1][]{
  colback=fondoPend, colframe=pink!60!black, boxrule=0.6pt,
  arc=2pt, left=6pt, right=6pt, top=5pt, bottom=5pt,
  title={Idea Random}, fonttitle=\bfseries\small, #1
}

% --- el resultado al que llego un ejercicio, para que se vea de
%     lejos al hojear la hoja antes del examen ---
\newtcolorbox{resultado}[1][]{
  colback=fondoResultado, colframe=green!45!black, boxrule=0.6pt,
  arc=2pt, left=6pt, right=6pt, top=5pt, bottom=5pt,
  title={Resultado}, fonttitle=\bfseries\small, #1
}

% --- código ---
\lstdefinestyle{codigo}{
  backgroundcolor=\color{fondoCodigo},
  basicstyle=\ttfamily\footnotesize,
  keywordstyle=\color{acento}\bfseries,
  commentstyle=\color{comentario}\itshape,
  stringstyle=\color{cadena},
  numbers=left, numberstyle=\tiny\color{gray}, numbersep=8pt,
  frame=single, rulecolor=\color{gray!40},
  breaklines=true, showstringspaces=false,
  tabsize=4, captionpos=b,
  extendedchars=true, inputencoding=utf8,
  literate={á}{{\'a}}1 {é}{{\'e}}1 {í}{{\'i}}1
           {ó}{{\'o}}1 {ú}{{\'u}}1 {ñ}{{\~n}}1
           {Ü}{{\"U}}1 {ü}{{\"u}}1 {¿}{{?`}}1
           {¡}{{!`}}1
}
\lstset{style=codigo}
\renewcommand{\lstlistingname}{Código}
\renewcommand{\lstlistlistingname}{Índice de código}
% --- pseudocódigo; algorithmicx arma el cuerpo y algorithm el
%     flotante que lo numera y lo pone en el índice. Las palabras
%     clave vienen en inglés de fábrica y babel no las toca, así que
%     se renombran una por una ---
\floatname{algorithm}{Algoritmo}
\renewcommand{\listalgorithmname}{Índice de algoritmos}

\algrenewcommand\algorithmicrequire{\textbf{Entrada:}}
\algrenewcommand\algorithmicensure{\textbf{Salida:}}
\algrenewcommand\algorithmicend{\textbf{fin}}
\algrenewcommand\algorithmicif{\textbf{si}}
\algrenewcommand\algorithmicthen{\textbf{entonces}}
\algrenewcommand\algorithmicelse{\textbf{si no}}
\algrenewcommand\algorithmicfor{\textbf{para}}
\algrenewcommand\algorithmicforall{\textbf{para todo}}
\algrenewcommand\algorithmicdo{\textbf{hacer}}
\algrenewcommand\algorithmicwhile{\textbf{mientras}}
\algrenewcommand\algorithmicrepeat{\textbf{repetir}}
\algrenewcommand\algorithmicuntil{\textbf{hasta que}}
\algrenewcommand\algorithmicloop{\textbf{ciclo}}
\algrenewcommand\algorithmicprocedure{\textbf{procedimiento}}
\algrenewcommand\algorithmicfunction{\textbf{función}}
\algrenewcommand\algorithmicreturn{\textbf{devolver}}
\algrenewcomment[1]{\hfill{\color{comentario}\itshape // #1}}

% Las dos que algorithmicx no trae. El \ del final es obligatorio: sin
% el TeX se come el espacio y sale "hastan-1" en vez de "hasta n-1"
\algnewcommand\Hasta{\textbf{hasta}\ }
\algnewcommand\Intercambiar{\textbf{intercambiar}\ }

% --- abreviaturas que se usan a cada rato ---
\newcommand{\N}{\mathbb{N}}
\newcommand{\Z}{\mathbb{Z}}
\newcommand{\Q}{\mathbb{Q}}
\newcommand{\R}{\mathbb{R}}
\newcommand{\C}{\mathbb{C}}
\newcommand{\abs}[1]{\left\lvert #1 \right\rvert}
\newcommand{\norm}[1]{\left\lVert #1 \right\rVert}
\newcommand{\set}[1]{\left\{ #1 \right\}}
\newcommand{\Ord}[1]{\mathcal{O}\!\left(#1\right)}
% Mayuscula la notacion grande, minuscula la chica
\newcommand{\Omg}[1]{\Omega\!\left(#1\right)}
\newcommand{\Tht}[1]{\Theta\!\left(#1\right)}
\newcommand{\ord}[1]{o\!\left(#1\right)}
\newcommand{\omg}[1]{\omega\!\left(#1\right)}
% Los tres limites sobre n, que salen en cada renglon del criterio
% del cociente para la notacion asintotica
\newcommand{\LM}{\lim_{n \to \infty}}
\newcommand{\LS}{\limsup_{n \to \infty}}
\newcommand{\LI}{\liminf_{n \to \infty}}

% --- encabezado de cada sesión: título a la izquierda, fecha a la
%     derecha. El argumento corto es el que sale en el índice ---
\newcommand{\clase}[2]{%
  \section[#1]{#1 \hfill \normalfont\small\textit{#2}}%
}

% --- portadilla y tabla de contenidos. Se llama una vez, justo
%     despues de \begin{document}. El renglon del profesor va dentro
%     del center, asi que \profesor puede traer varios nombres
%     separados por \\ y salen apilados y centrados sin tocar nada ---
\newcommand{\portadilla}{%
  \begin{center}
    {\LARGE\bfseries\materia\par}
    \vspace{0.4em}
    {\large \profesor \par}
    \vspace{0.2em}
    {\small \periodo\quad-\quad\alumno\par}
    \vspace{0.3em}
    \rule{\linewidth}{0.5pt}
  \end{center}
  \vspace{0.5em}
  \tableofcontents
  \vspace{1em}
  \rule{\linewidth}{0.5pt}%
}

% =====================================================================
%  DOCUMENTOS DE SOLUCIONES
%  Lo que usa una tanda de ejercicios resueltos. Una libreta no toca
%  nada de aqui, pero vive en el preambulo compartido para no tener
%  dos archivos que arreglar cuando cambie el estilo.
% =====================================================================

% --- encabezado compacto: titulo, subtitulo y regla. Sustituye a
%     \portadilla cuando el documento es una tanda de ejercicios: sin
%     indice, porque se hojea de corrido, no se navega ---
\newcommand{\encabezadoSoluciones}[2]{%
  % El paquete algorithm se carga con la opcion [section] para que en
  % una libreta el pseudocodigo se numere por clase. Aqui no hay
  % secciones, asi que esa numeracion sale como "Algoritmo 0.1". Se
  % corrige al vuelo, y solo en los documentos que llaman a este
  % encabezado: la libreta no se entera.
  \renewcommand{\thealgorithm}{\arabic{algorithm}}%
  \begin{center}
    {\LARGE\bfseries #1\par}
    \vspace{0.3em}
    {\small\itshape #2\par}
    \vspace{0.5em}
    \rule{\linewidth}{0.5pt}
  \end{center}
  \vspace{0.3em}%
}

% --- hoja de un solo ejercicio: un renglon chico arriba con el
%     titulo y su numero, y ya. Sin titulo grande y sin pie de pagina,
%     porque en una hoja de dos paginas cada centimetro que se va en
%     adornos es un paso del desarrollo que no cabe ---
\newcommand{\encabezadoEjercicio}[3]{%
  % Misma correccion que en \encabezadoSoluciones: sin secciones, el
  % pseudocodigo saldria numerado "0.1"
  \renewcommand{\thealgorithm}{\arabic{algorithm}}%
  \fancyhf{}%
  \fancyhead[L]{\small\textsc{#1} (ejercicio #2 de #3)}%
  % El nombre va a la derecha porque las hojas de respuestas se
  % entregan y el profesor pide que cada una lo lleve
  \fancyhead[R]{\small\alumno}%
  \renewcommand{\headrulewidth}{0.4pt}%
  \pagestyle{fancy}%
  \thispagestyle{fancy}%
}

% --- bloque tematico, numerado con romanos: I, II, III. Agrupa los
%     ejercicios que comparten tema. No dibuja regla: la del
%     \problema que sigue ya separa el titulo del primer ejercicio,
%     y dos rayas seguidas dejan un hueco feo ---
\newcounter{parte}
\renewcommand{\theparte}{\Roman{parte}}
\newcommand{\parte}[1]{%
  \bigskip
  \refstepcounter{parte}%
  \noindent{\large\bfseries \theparte. #1\par}%
}

% --- un ejercicio resuelto. Se numera solo y corrido a lo largo del
%     documento, sin reiniciar en cada parte, para poder decir "el 14"
%     sin aclarar de que bloque ---
\newcounter{problema}
\newcommand{\problema}[1]{%
  \bigskip
  \noindent\rule{\linewidth}{0.4pt}\par
  \vspace{0.3em}
  \refstepcounter{problema}%
  \noindent{\bfseries Ejercicio \theproblema\ --- #1\par}
  \vspace{0.3em}%
}

% --- rotulo del paso que sigue: Caso base, Hipotesis inductiva, Paso
%     inductivo. El texto sigue en el mismo renglon ---
\newcommand{\paso}[1]{\par\vspace{0.3em}\noindent\textbf{#1}\ }

% --- la respuesta de un ejercicio, centrada y en un cuadro. El
%     argumento va en modo matematico, sin los delimitadores ---
\newcommand{\respuesta}[1]{%
  \[
    \boxed{\;#1\;}
  \]
}

% --- resaltado inline, como pasarle el plumon a media frase. No
%     admite matematicas adentro: soul rompe con $...$, asi que una
%     formula que hay que destacar va en \respuesta ---
\sethlcolor{marcador}
\newcommand{\resaltar}[1]{\hl{#1}}

% --- justificacion al final de un renglon de \align, en chiquito y
%     entre parentesis, para decir de donde salio el paso ---
\newcommand{\razon}[1]{\quad\text{\small\itshape (#1)}}

%############################ END OF PREAMBULO.TEX #############################
%###############################################################################

%  \begin{document}
%  \portadilla
%
%  Con dos o mas profesores se declara UN solo \profesor y los nombres
%  se separan con \\, uno por renglon:
%
%  \newcommand{\profesor}{Dra. Fulana de Tal \\
%                         Dr. Mengano de Cual}
% =====================================================================


% ==================== IDIOMA Y CODIFICACIÓN ====================
\usepackage[utf8]{inputenc}
\usepackage[T1]{fontenc}
\usepackage[spanish,es-tabla,es-noshorthands]{babel}
\usepackage{lmodern}
\usepackage{microtype}

% ==================== PÁGINA ====================
\usepackage[letterpaper,margin=2.3cm,headheight=14pt]{geometry}
\usepackage{parskip}
\usepackage{fancyhdr}
\usepackage{enumitem}

% ==================== MATEMÁTICAS ====================
\usepackage{amsmath,amssymb,amsthm,mathtools}

% ==================== FIGURAS, TABLAS, CÓDIGO ====================
\usepackage{graphicx}
\usepackage{booktabs}
\usepackage{array}
% float da el especificador [H], que clava una tabla donde se
% escribio. Una traza que LaTeX manda dos paginas adelante deja
% de leerse junto al paso que la produjo
\usepackage{float}
\usepackage{xcolor}
% tikz ya viene arrastrado por tcolorbox[most], pero se carga
% explicito para no depender de ese detalle; la libreria babel
% protege a tikz de los caracteres activos del español
\usepackage{tikz}
\usetikzlibrary{babel}
% pgfplots dibuja graficas de datos (ejes, escalas log) sobre tikz
\usepackage{pgfplots}
\pgfplotsset{compat=1.18}
\usepackage{listings}
\usepackage[ruled,section]{algorithm}
\usepackage{algpseudocode}
\usepackage[most]{tcolorbox}
% soul da \hl, el unico resaltado que parte renglon; \colorbox no
\usepackage{soul}
\usepackage{hyperref}

% =====================================================================
%  ESTILO
% =====================================================================

% --- encabezado y pie de cada página ---
\pagestyle{fancy}
\fancyhf{}
\fancyhead[L]{\small\textsc{\materia}}
\fancyhead[R]{\small\periodo}
\fancyfoot[C]{\small\thepage}
\renewcommand{\headrulewidth}{0.4pt}

% --- enlaces ---
\hypersetup{
  colorlinks=true,
  linkcolor=black,
  urlcolor=blue,
  citecolor=black,
  pdftitle={Apuntes de \materia},
  pdfauthor={\alumno}
}

% --- paleta ---
\definecolor{acento}{HTML}{1F4E79}
\definecolor{fondoIdea}{HTML}{EAF2FB}
\definecolor{fondoOjo}{HTML}{FDF0E6}
\definecolor{fondoPend}{HTML}{F2F0F7}
\definecolor{comentario}{HTML}{5A7D5A}
\definecolor{cadena}{HTML}{A03030}
\definecolor{fondoCodigo}{HTML}{F7F7F7}
\definecolor{fondoResultado}{HTML}{EAF5EC}
\definecolor{marcador}{HTML}{FFF3A3}

% --- entornos de teorema; numeran por sección, o sea por clase ---
\theoremstyle{definition}
\newtheorem{definicion}{Definición}[section]
\newtheorem{ejemplo}[definicion]{Ejemplo}
\newtheorem{ejercicio}[definicion]{Ejercicio}

\theoremstyle{plain}
\newtheorem{teorema}[definicion]{Teorema}
\newtheorem{lema}[definicion]{Lema}
\newtheorem{corolario}[definicion]{Corolario}
\newtheorem{proposicion}[definicion]{Proposición}

\theoremstyle{remark}
\newtheorem*{observacion}{Observación}
\newtheorem*{quick_sum}{Resumen rápido}

% --- cajas de color para lo que se relee antes del examen ---
\newtcolorbox{idea}[1][]{
  colback=fondoIdea, colframe=acento, boxrule=0.6pt,
  arc=2pt, left=6pt, right=6pt, top=5pt, bottom=5pt,
  title={Idea clave}, fonttitle=\bfseries\small, #1
}
\newtcolorbox{ojo}[1][]{
  colback=fondoOjo, colframe=orange!70!black, boxrule=0.6pt,
  arc=2pt, left=6pt, right=6pt, top=5pt, bottom=5pt,
  title={Ojito...}, fonttitle=\bfseries\small, #1
}
\newtcolorbox{pendiente}[1][]{
  colback=fondoPend, colframe=violet!60!black, boxrule=0.6pt,
  arc=2pt, left=6pt, right=6pt, top=5pt, bottom=5pt,
  title={Pendiente por resolver}, fonttitle=\bfseries\small, #1
}
\newtcolorbox{idear}[1][]{
  colback=fondoPend, colframe=pink!60!black, boxrule=0.6pt,
  arc=2pt, left=6pt, right=6pt, top=5pt, bottom=5pt,
  title={Idea Random}, fonttitle=\bfseries\small, #1
}

% --- el resultado al que llego un ejercicio, para que se vea de
%     lejos al hojear la hoja antes del examen ---
\newtcolorbox{resultado}[1][]{
  colback=fondoResultado, colframe=green!45!black, boxrule=0.6pt,
  arc=2pt, left=6pt, right=6pt, top=5pt, bottom=5pt,
  title={Resultado}, fonttitle=\bfseries\small, #1
}

% --- código ---
\lstdefinestyle{codigo}{
  backgroundcolor=\color{fondoCodigo},
  basicstyle=\ttfamily\footnotesize,
  keywordstyle=\color{acento}\bfseries,
  commentstyle=\color{comentario}\itshape,
  stringstyle=\color{cadena},
  numbers=left, numberstyle=\tiny\color{gray}, numbersep=8pt,
  frame=single, rulecolor=\color{gray!40},
  breaklines=true, showstringspaces=false,
  tabsize=4, captionpos=b,
  extendedchars=true, inputencoding=utf8,
  literate={á}{{\'a}}1 {é}{{\'e}}1 {í}{{\'i}}1
           {ó}{{\'o}}1 {ú}{{\'u}}1 {ñ}{{\~n}}1
           {Ü}{{\"U}}1 {ü}{{\"u}}1 {¿}{{?`}}1
           {¡}{{!`}}1
}
\lstset{style=codigo}
\renewcommand{\lstlistingname}{Código}
\renewcommand{\lstlistlistingname}{Índice de código}
% --- pseudocódigo; algorithmicx arma el cuerpo y algorithm el
%     flotante que lo numera y lo pone en el índice. Las palabras
%     clave vienen en inglés de fábrica y babel no las toca, así que
%     se renombran una por una ---
\floatname{algorithm}{Algoritmo}
\renewcommand{\listalgorithmname}{Índice de algoritmos}

\algrenewcommand\algorithmicrequire{\textbf{Entrada:}}
\algrenewcommand\algorithmicensure{\textbf{Salida:}}
\algrenewcommand\algorithmicend{\textbf{fin}}
\algrenewcommand\algorithmicif{\textbf{si}}
\algrenewcommand\algorithmicthen{\textbf{entonces}}
\algrenewcommand\algorithmicelse{\textbf{si no}}
\algrenewcommand\algorithmicfor{\textbf{para}}
\algrenewcommand\algorithmicforall{\textbf{para todo}}
\algrenewcommand\algorithmicdo{\textbf{hacer}}
\algrenewcommand\algorithmicwhile{\textbf{mientras}}
\algrenewcommand\algorithmicrepeat{\textbf{repetir}}
\algrenewcommand\algorithmicuntil{\textbf{hasta que}}
\algrenewcommand\algorithmicloop{\textbf{ciclo}}
\algrenewcommand\algorithmicprocedure{\textbf{procedimiento}}
\algrenewcommand\algorithmicfunction{\textbf{función}}
\algrenewcommand\algorithmicreturn{\textbf{devolver}}
\algrenewcomment[1]{\hfill{\color{comentario}\itshape // #1}}

% Las dos que algorithmicx no trae. El \ del final es obligatorio: sin
% el TeX se come el espacio y sale "hastan-1" en vez de "hasta n-1"
\algnewcommand\Hasta{\textbf{hasta}\ }
\algnewcommand\Intercambiar{\textbf{intercambiar}\ }

% --- abreviaturas que se usan a cada rato ---
\newcommand{\N}{\mathbb{N}}
\newcommand{\Z}{\mathbb{Z}}
\newcommand{\Q}{\mathbb{Q}}
\newcommand{\R}{\mathbb{R}}
\newcommand{\C}{\mathbb{C}}
\newcommand{\abs}[1]{\left\lvert #1 \right\rvert}
\newcommand{\norm}[1]{\left\lVert #1 \right\rVert}
\newcommand{\set}[1]{\left\{ #1 \right\}}
\newcommand{\Ord}[1]{\mathcal{O}\!\left(#1\right)}
% Mayuscula la notacion grande, minuscula la chica
\newcommand{\Omg}[1]{\Omega\!\left(#1\right)}
\newcommand{\Tht}[1]{\Theta\!\left(#1\right)}
\newcommand{\ord}[1]{o\!\left(#1\right)}
\newcommand{\omg}[1]{\omega\!\left(#1\right)}
% Los tres limites sobre n, que salen en cada renglon del criterio
% del cociente para la notacion asintotica
\newcommand{\LM}{\lim_{n \to \infty}}
\newcommand{\LS}{\limsup_{n \to \infty}}
\newcommand{\LI}{\liminf_{n \to \infty}}

% --- encabezado de cada sesión: título a la izquierda, fecha a la
%     derecha. El argumento corto es el que sale en el índice ---
\newcommand{\clase}[2]{%
  \section[#1]{#1 \hfill \normalfont\small\textit{#2}}%
}

% --- portadilla y tabla de contenidos. Se llama una vez, justo
%     despues de \begin{document}. El renglon del profesor va dentro
%     del center, asi que \profesor puede traer varios nombres
%     separados por \\ y salen apilados y centrados sin tocar nada ---
\newcommand{\portadilla}{%
  \begin{center}
    {\LARGE\bfseries\materia\par}
    \vspace{0.4em}
    {\large \profesor \par}
    \vspace{0.2em}
    {\small \periodo\quad-\quad\alumno\par}
    \vspace{0.3em}
    \rule{\linewidth}{0.5pt}
  \end{center}
  \vspace{0.5em}
  \tableofcontents
  \vspace{1em}
  \rule{\linewidth}{0.5pt}%
}

% =====================================================================
%  DOCUMENTOS DE SOLUCIONES
%  Lo que usa una tanda de ejercicios resueltos. Una libreta no toca
%  nada de aqui, pero vive en el preambulo compartido para no tener
%  dos archivos que arreglar cuando cambie el estilo.
% =====================================================================

% --- encabezado compacto: titulo, subtitulo y regla. Sustituye a
%     \portadilla cuando el documento es una tanda de ejercicios: sin
%     indice, porque se hojea de corrido, no se navega ---
\newcommand{\encabezadoSoluciones}[2]{%
  % El paquete algorithm se carga con la opcion [section] para que en
  % una libreta el pseudocodigo se numere por clase. Aqui no hay
  % secciones, asi que esa numeracion sale como "Algoritmo 0.1". Se
  % corrige al vuelo, y solo en los documentos que llaman a este
  % encabezado: la libreta no se entera.
  \renewcommand{\thealgorithm}{\arabic{algorithm}}%
  \begin{center}
    {\LARGE\bfseries #1\par}
    \vspace{0.3em}
    {\small\itshape #2\par}
    \vspace{0.5em}
    \rule{\linewidth}{0.5pt}
  \end{center}
  \vspace{0.3em}%
}

% --- hoja de un solo ejercicio: un renglon chico arriba con el
%     titulo y su numero, y ya. Sin titulo grande y sin pie de pagina,
%     porque en una hoja de dos paginas cada centimetro que se va en
%     adornos es un paso del desarrollo que no cabe ---
\newcommand{\encabezadoEjercicio}[3]{%
  % Misma correccion que en \encabezadoSoluciones: sin secciones, el
  % pseudocodigo saldria numerado "0.1"
  \renewcommand{\thealgorithm}{\arabic{algorithm}}%
  \fancyhf{}%
  \fancyhead[L]{\small\textsc{#1} (ejercicio #2 de #3)}%
  % El nombre va a la derecha porque las hojas de respuestas se
  % entregan y el profesor pide que cada una lo lleve
  \fancyhead[R]{\small\alumno}%
  \renewcommand{\headrulewidth}{0.4pt}%
  \pagestyle{fancy}%
  \thispagestyle{fancy}%
}

% --- bloque tematico, numerado con romanos: I, II, III. Agrupa los
%     ejercicios que comparten tema. No dibuja regla: la del
%     \problema que sigue ya separa el titulo del primer ejercicio,
%     y dos rayas seguidas dejan un hueco feo ---
\newcounter{parte}
\renewcommand{\theparte}{\Roman{parte}}
\newcommand{\parte}[1]{%
  \bigskip
  \refstepcounter{parte}%
  \noindent{\large\bfseries \theparte. #1\par}%
}

% --- un ejercicio resuelto. Se numera solo y corrido a lo largo del
%     documento, sin reiniciar en cada parte, para poder decir "el 14"
%     sin aclarar de que bloque ---
\newcounter{problema}
\newcommand{\problema}[1]{%
  \bigskip
  \noindent\rule{\linewidth}{0.4pt}\par
  \vspace{0.3em}
  \refstepcounter{problema}%
  \noindent{\bfseries Ejercicio \theproblema\ --- #1\par}
  \vspace{0.3em}%
}

% --- rotulo del paso que sigue: Caso base, Hipotesis inductiva, Paso
%     inductivo. El texto sigue en el mismo renglon ---
\newcommand{\paso}[1]{\par\vspace{0.3em}\noindent\textbf{#1}\ }

% --- la respuesta de un ejercicio, centrada y en un cuadro. El
%     argumento va en modo matematico, sin los delimitadores ---
\newcommand{\respuesta}[1]{%
  \[
    \boxed{\;#1\;}
  \]
}

% --- resaltado inline, como pasarle el plumon a media frase. No
%     admite matematicas adentro: soul rompe con $...$, asi que una
%     formula que hay que destacar va en \respuesta ---
\sethlcolor{marcador}
\newcommand{\resaltar}[1]{\hl{#1}}

% --- justificacion al final de un renglon de \align, en chiquito y
%     entre parentesis, para decir de donde salio el paso ---
\newcommand{\razon}[1]{\quad\text{\small\itshape (#1)}}

%############################ END OF PREAMBULO.TEX #############################
%###############################################################################

%  \begin{document}
%  \portadilla
%
%  Con dos o mas profesores se declara UN solo \profesor y los nombres
%  se separan con \\, uno por renglon:
%
%  \newcommand{\profesor}{Dra. Fulana de Tal \\
%                         Dr. Mengano de Cual}
% =====================================================================


% ==================== IDIOMA Y CODIFICACIÓN ====================
\usepackage[utf8]{inputenc}
\usepackage[T1]{fontenc}
\usepackage[spanish,es-tabla,es-noshorthands]{babel}
\usepackage{lmodern}
\usepackage{microtype}

% ==================== PÁGINA ====================
\usepackage[letterpaper,margin=2.3cm,headheight=14pt]{geometry}
\usepackage{parskip}
\usepackage{fancyhdr}
\usepackage{enumitem}

% ==================== MATEMÁTICAS ====================
\usepackage{amsmath,amssymb,amsthm,mathtools}

% ==================== FIGURAS, TABLAS, CÓDIGO ====================
\usepackage{graphicx}
\usepackage{booktabs}
\usepackage{array}
% float da el especificador [H], que clava una tabla donde se
% escribio. Una traza que LaTeX manda dos paginas adelante deja
% de leerse junto al paso que la produjo
\usepackage{float}
\usepackage{xcolor}
% tikz ya viene arrastrado por tcolorbox[most], pero se carga
% explicito para no depender de ese detalle; la libreria babel
% protege a tikz de los caracteres activos del español
\usepackage{tikz}
\usetikzlibrary{babel}
% pgfplots dibuja graficas de datos (ejes, escalas log) sobre tikz
\usepackage{pgfplots}
\pgfplotsset{compat=1.18}
\usepackage{listings}
\usepackage[ruled,section]{algorithm}
\usepackage{algpseudocode}
\usepackage[most]{tcolorbox}
% soul da \hl, el unico resaltado que parte renglon; \colorbox no
\usepackage{soul}
\usepackage{hyperref}

% =====================================================================
%  ESTILO
% =====================================================================

% --- encabezado y pie de cada página ---
\pagestyle{fancy}
\fancyhf{}
\fancyhead[L]{\small\textsc{\materia}}
\fancyhead[R]{\small\periodo}
\fancyfoot[C]{\small\thepage}
\renewcommand{\headrulewidth}{0.4pt}

% --- enlaces ---
\hypersetup{
  colorlinks=true,
  linkcolor=black,
  urlcolor=blue,
  citecolor=black,
  pdftitle={Apuntes de \materia},
  pdfauthor={\alumno}
}

% --- paleta ---
\definecolor{acento}{HTML}{1F4E79}
\definecolor{fondoIdea}{HTML}{EAF2FB}
\definecolor{fondoOjo}{HTML}{FDF0E6}
\definecolor{fondoPend}{HTML}{F2F0F7}
\definecolor{comentario}{HTML}{5A7D5A}
\definecolor{cadena}{HTML}{A03030}
\definecolor{fondoCodigo}{HTML}{F7F7F7}
\definecolor{fondoResultado}{HTML}{EAF5EC}
\definecolor{marcador}{HTML}{FFF3A3}

% --- entornos de teorema; numeran por sección, o sea por clase ---
\theoremstyle{definition}
\newtheorem{definicion}{Definición}[section]
\newtheorem{ejemplo}[definicion]{Ejemplo}
\newtheorem{ejercicio}[definicion]{Ejercicio}

\theoremstyle{plain}
\newtheorem{teorema}[definicion]{Teorema}
\newtheorem{lema}[definicion]{Lema}
\newtheorem{corolario}[definicion]{Corolario}
\newtheorem{proposicion}[definicion]{Proposición}

\theoremstyle{remark}
\newtheorem*{observacion}{Observación}
\newtheorem*{quick_sum}{Resumen rápido}

% --- cajas de color para lo que se relee antes del examen ---
\newtcolorbox{idea}[1][]{
  colback=fondoIdea, colframe=acento, boxrule=0.6pt,
  arc=2pt, left=6pt, right=6pt, top=5pt, bottom=5pt,
  title={Idea clave}, fonttitle=\bfseries\small, #1
}
\newtcolorbox{ojo}[1][]{
  colback=fondoOjo, colframe=orange!70!black, boxrule=0.6pt,
  arc=2pt, left=6pt, right=6pt, top=5pt, bottom=5pt,
  title={Ojito...}, fonttitle=\bfseries\small, #1
}
\newtcolorbox{pendiente}[1][]{
  colback=fondoPend, colframe=violet!60!black, boxrule=0.6pt,
  arc=2pt, left=6pt, right=6pt, top=5pt, bottom=5pt,
  title={Pendiente por resolver}, fonttitle=\bfseries\small, #1
}
\newtcolorbox{idear}[1][]{
  colback=fondoPend, colframe=pink!60!black, boxrule=0.6pt,
  arc=2pt, left=6pt, right=6pt, top=5pt, bottom=5pt,
  title={Idea Random}, fonttitle=\bfseries\small, #1
}

% --- el resultado al que llego un ejercicio, para que se vea de
%     lejos al hojear la hoja antes del examen ---
\newtcolorbox{resultado}[1][]{
  colback=fondoResultado, colframe=green!45!black, boxrule=0.6pt,
  arc=2pt, left=6pt, right=6pt, top=5pt, bottom=5pt,
  title={Resultado}, fonttitle=\bfseries\small, #1
}

% --- código ---
\lstdefinestyle{codigo}{
  backgroundcolor=\color{fondoCodigo},
  basicstyle=\ttfamily\footnotesize,
  keywordstyle=\color{acento}\bfseries,
  commentstyle=\color{comentario}\itshape,
  stringstyle=\color{cadena},
  numbers=left, numberstyle=\tiny\color{gray}, numbersep=8pt,
  frame=single, rulecolor=\color{gray!40},
  breaklines=true, showstringspaces=false,
  tabsize=4, captionpos=b,
  extendedchars=true, inputencoding=utf8,
  literate={á}{{\'a}}1 {é}{{\'e}}1 {í}{{\'i}}1
           {ó}{{\'o}}1 {ú}{{\'u}}1 {ñ}{{\~n}}1
           {Ü}{{\"U}}1 {ü}{{\"u}}1 {¿}{{?`}}1
           {¡}{{!`}}1
}
\lstset{style=codigo}
\renewcommand{\lstlistingname}{Código}
\renewcommand{\lstlistlistingname}{Índice de código}
% --- pseudocódigo; algorithmicx arma el cuerpo y algorithm el
%     flotante que lo numera y lo pone en el índice. Las palabras
%     clave vienen en inglés de fábrica y babel no las toca, así que
%     se renombran una por una ---
\floatname{algorithm}{Algoritmo}
\renewcommand{\listalgorithmname}{Índice de algoritmos}

\algrenewcommand\algorithmicrequire{\textbf{Entrada:}}
\algrenewcommand\algorithmicensure{\textbf{Salida:}}
\algrenewcommand\algorithmicend{\textbf{fin}}
\algrenewcommand\algorithmicif{\textbf{si}}
\algrenewcommand\algorithmicthen{\textbf{entonces}}
\algrenewcommand\algorithmicelse{\textbf{si no}}
\algrenewcommand\algorithmicfor{\textbf{para}}
\algrenewcommand\algorithmicforall{\textbf{para todo}}
\algrenewcommand\algorithmicdo{\textbf{hacer}}
\algrenewcommand\algorithmicwhile{\textbf{mientras}}
\algrenewcommand\algorithmicrepeat{\textbf{repetir}}
\algrenewcommand\algorithmicuntil{\textbf{hasta que}}
\algrenewcommand\algorithmicloop{\textbf{ciclo}}
\algrenewcommand\algorithmicprocedure{\textbf{procedimiento}}
\algrenewcommand\algorithmicfunction{\textbf{función}}
\algrenewcommand\algorithmicreturn{\textbf{devolver}}
\algrenewcomment[1]{\hfill{\color{comentario}\itshape // #1}}

% Las dos que algorithmicx no trae. El \ del final es obligatorio: sin
% el TeX se come el espacio y sale "hastan-1" en vez de "hasta n-1"
\algnewcommand\Hasta{\textbf{hasta}\ }
\algnewcommand\Intercambiar{\textbf{intercambiar}\ }

% --- abreviaturas que se usan a cada rato ---
\newcommand{\N}{\mathbb{N}}
\newcommand{\Z}{\mathbb{Z}}
\newcommand{\Q}{\mathbb{Q}}
\newcommand{\R}{\mathbb{R}}
\newcommand{\C}{\mathbb{C}}
\newcommand{\abs}[1]{\left\lvert #1 \right\rvert}
\newcommand{\norm}[1]{\left\lVert #1 \right\rVert}
\newcommand{\set}[1]{\left\{ #1 \right\}}
\newcommand{\Ord}[1]{\mathcal{O}\!\left(#1\right)}
% Mayuscula la notacion grande, minuscula la chica
\newcommand{\Omg}[1]{\Omega\!\left(#1\right)}
\newcommand{\Tht}[1]{\Theta\!\left(#1\right)}
\newcommand{\ord}[1]{o\!\left(#1\right)}
\newcommand{\omg}[1]{\omega\!\left(#1\right)}
% Los tres limites sobre n, que salen en cada renglon del criterio
% del cociente para la notacion asintotica
\newcommand{\LM}{\lim_{n \to \infty}}
\newcommand{\LS}{\limsup_{n \to \infty}}
\newcommand{\LI}{\liminf_{n \to \infty}}

% --- encabezado de cada sesión: título a la izquierda, fecha a la
%     derecha. El argumento corto es el que sale en el índice ---
\newcommand{\clase}[2]{%
  \section[#1]{#1 \hfill \normalfont\small\textit{#2}}%
}

% --- portadilla y tabla de contenidos. Se llama una vez, justo
%     despues de \begin{document}. El renglon del profesor va dentro
%     del center, asi que \profesor puede traer varios nombres
%     separados por \\ y salen apilados y centrados sin tocar nada ---
\newcommand{\portadilla}{%
  \begin{center}
    {\LARGE\bfseries\materia\par}
    \vspace{0.4em}
    {\large \profesor \par}
    \vspace{0.2em}
    {\small \periodo\quad-\quad\alumno\par}
    \vspace{0.3em}
    \rule{\linewidth}{0.5pt}
  \end{center}
  \vspace{0.5em}
  \tableofcontents
  \vspace{1em}
  \rule{\linewidth}{0.5pt}%
}

% =====================================================================
%  DOCUMENTOS DE SOLUCIONES
%  Lo que usa una tanda de ejercicios resueltos. Una libreta no toca
%  nada de aqui, pero vive en el preambulo compartido para no tener
%  dos archivos que arreglar cuando cambie el estilo.
% =====================================================================

% --- encabezado compacto: titulo, subtitulo y regla. Sustituye a
%     \portadilla cuando el documento es una tanda de ejercicios: sin
%     indice, porque se hojea de corrido, no se navega ---
\newcommand{\encabezadoSoluciones}[2]{%
  % El paquete algorithm se carga con la opcion [section] para que en
  % una libreta el pseudocodigo se numere por clase. Aqui no hay
  % secciones, asi que esa numeracion sale como "Algoritmo 0.1". Se
  % corrige al vuelo, y solo en los documentos que llaman a este
  % encabezado: la libreta no se entera.
  \renewcommand{\thealgorithm}{\arabic{algorithm}}%
  \begin{center}
    {\LARGE\bfseries #1\par}
    \vspace{0.3em}
    {\small\itshape #2\par}
    \vspace{0.5em}
    \rule{\linewidth}{0.5pt}
  \end{center}
  \vspace{0.3em}%
}

% --- hoja de un solo ejercicio: un renglon chico arriba con el
%     titulo y su numero, y ya. Sin titulo grande y sin pie de pagina,
%     porque en una hoja de dos paginas cada centimetro que se va en
%     adornos es un paso del desarrollo que no cabe ---
\newcommand{\encabezadoEjercicio}[3]{%
  % Misma correccion que en \encabezadoSoluciones: sin secciones, el
  % pseudocodigo saldria numerado "0.1"
  \renewcommand{\thealgorithm}{\arabic{algorithm}}%
  \fancyhf{}%
  \fancyhead[L]{\small\textsc{#1} (ejercicio #2 de #3)}%
  % El nombre va a la derecha porque las hojas de respuestas se
  % entregan y el profesor pide que cada una lo lleve
  \fancyhead[R]{\small\alumno}%
  \renewcommand{\headrulewidth}{0.4pt}%
  \pagestyle{fancy}%
  \thispagestyle{fancy}%
}

% --- bloque tematico, numerado con romanos: I, II, III. Agrupa los
%     ejercicios que comparten tema. No dibuja regla: la del
%     \problema que sigue ya separa el titulo del primer ejercicio,
%     y dos rayas seguidas dejan un hueco feo ---
\newcounter{parte}
\renewcommand{\theparte}{\Roman{parte}}
\newcommand{\parte}[1]{%
  \bigskip
  \refstepcounter{parte}%
  \noindent{\large\bfseries \theparte. #1\par}%
}

% --- un ejercicio resuelto. Se numera solo y corrido a lo largo del
%     documento, sin reiniciar en cada parte, para poder decir "el 14"
%     sin aclarar de que bloque ---
\newcounter{problema}
\newcommand{\problema}[1]{%
  \bigskip
  \noindent\rule{\linewidth}{0.4pt}\par
  \vspace{0.3em}
  \refstepcounter{problema}%
  \noindent{\bfseries Ejercicio \theproblema\ --- #1\par}
  \vspace{0.3em}%
}

% --- rotulo del paso que sigue: Caso base, Hipotesis inductiva, Paso
%     inductivo. El texto sigue en el mismo renglon ---
\newcommand{\paso}[1]{\par\vspace{0.3em}\noindent\textbf{#1}\ }

% --- la respuesta de un ejercicio, centrada y en un cuadro. El
%     argumento va en modo matematico, sin los delimitadores ---
\newcommand{\respuesta}[1]{%
  \[
    \boxed{\;#1\;}
  \]
}

% --- resaltado inline, como pasarle el plumon a media frase. No
%     admite matematicas adentro: soul rompe con $...$, asi que una
%     formula que hay que destacar va en \respuesta ---
\sethlcolor{marcador}
\newcommand{\resaltar}[1]{\hl{#1}}

% --- justificacion al final de un renglon de \align, en chiquito y
%     entre parentesis, para decir de donde salio el paso ---
\newcommand{\razon}[1]{\quad\text{\small\itshape (#1)}}

%############################ END OF PREAMBULO.TEX #############################
%###############################################################################


% =====================================================================
%  NOTACIÓN PROPIA DE ESTE DOCUMENTO
% =====================================================================
\DeclareMathOperator*{\argmax}{arg\,max}
\newcommand{\Count}{\mathrm{C}}
\newcommand{\Vocab}{\mathcal{V}}
\newcommand{\bos}{\texttt{<s>}}
\newcommand{\eos}{\texttt{</s>}}
\newcommand{\unk}{\texttt{<UNK>}}

\title{\textbf{Modelos de lenguaje de $n$-gramas}\\[2pt]
  \large Factorización, supuesto de Markov, estimación, suavizado y
  perplejidad (ref.\ J\&M SLP \S3)}
\author{Material de apoyo para \materia, INAOE}
\date{}

% =====================================================================
\begin{document}
\maketitle

\section{Qué objeto es un modelo de lenguaje}

Antes de escribir cualquier fórmula conviene fijar qué se está modelando, porque
la mayor parte de las confusiones posteriores vienen de perder de vista el
objeto primario.

\begin{definicion}
Un \emph{modelo de lenguaje} (language model) es una distribución de
probabilidad sobre $\Vocab^*$, el conjunto de todas las cadenas finitas sobre un
vocabulario finito $\Vocab$. Asigna a cada secuencia $w_1\cdots w_N$ un número
$P(w_1\cdots w_N)\ge 0$ sujeto a
\[
\sum_{s\in\Vocab^*} P(s) = 1 .
\]
\end{definicion}

\begin{idea}[title={Lo que hay que retener de entrada}]
El objeto primario es la \textbf{conjunta sobre secuencias completas}, no las
probabilidades condicionales. Las condicionales aparecerán enseguida, pero como
una \emph{manera de parametrizar} esa conjunta, no como el objeto que se define.
Esta distinción es la que se pierde cuando se afirma que ``el modelo está
normalizado'' sin decir sobre qué espacio.
\end{idea}

Nótese que $\Vocab^*$ es infinito numerable: contiene cadenas de toda longitud
finita. Una distribución sobre un conjunto infinito puede perfectamente sumar
menos que $1$ si parte de la masa se escapa hacia objetos que no pertenecen al
conjunto; a este fenómeno volveremos en la sección \ref{sec:fuga}, porque no es
una patología teórica sino algo que un modelo mal construido produce con
facilidad.

\section{La regla de la cadena: factorización exacta}

\subsection{Enunciado}

La regla de la cadena de la probabilidad (chain rule) no es más que la
definición de probabilidad condicional aplicada repetidamente. Para dos eventos
con $P(B)>0$,
\[
P(A\mid B)=\frac{P(A\cap B)}{P(B)}
\quad\Longrightarrow\quad
P(A\cap B) = P(B)\,P(A\mid B),
\]
y al iterar sobre $N$ eventos se obtiene
\[
P(A_1\cap\cdots\cap A_N)
= P(A_1)\,P(A_2\mid A_1)\,P(A_3\mid A_1\cap A_2)\cdots P(A_N\mid
A_1\cap\cdots\cap A_{N-1}).
\]

En el caso del modelo de lenguaje, $A_i$ es el evento ``la $i$-ésima palabra es
$w_i$'', y la intersección de todos ellos es el evento ``la secuencia es
$w_1\cdots w_N$''. Se obtiene entonces
\begin{equation}\label{eq:chain}
P(w_1\cdots w_N) \;=\; \prod_{i=1}^{N} P(w_i \mid w_1\cdots w_{i-1}).
\end{equation}

\begin{idea}[title={La factorización no aproxima nada}]
La ecuación \eqref{eq:chain} es una identidad, no un modelo. No se ha
introducido ningún supuesto sobre el lenguaje. El único problema que queda es de
\emph{estimación}: el condicionante $w_1\cdots w_{i-1}$ crece sin cota, de modo
que habría un número no acotado de distribuciones distintas que estimar, una por
cada historial posible.
\end{idea}

\subsection{Verificación por telescopio y por inducción}

Conviene ver la cancelación explícitamente antes de aceptar la fórmula general.
Para $N=3$, reescribiendo cada factor con la definición de condicional:
\[
P(w_1)\cdot\frac{P(w_1w_2)}{P(w_1)}\cdot\frac{P(w_1w_2w_3)}{P(w_1w_2)} \;=\;
P(w_1w_2w_3),
\]
donde $P(w_1)$ cancela con el denominador del segundo factor y $P(w_1w_2)$ con
el del tercero.

\begin{proposicion}
La identidad \eqref{eq:chain} vale para todo $N\ge 1$ siempre que $P(w_1\cdots
w_{i-1})>0$ para cada $i\le N$.
\end{proposicion}

\begin{proof}
Inducción sobre $N$, la longitud de la secuencia.

\emph{Base} ($N=1$): ambos lados son $P(w_1)$, tomando como convención que el
producto sobre el historial vacío vale $1$.

\emph{Paso inductivo}: supóngase la identidad válida para $N-1$. Aplíquese la
regla de dos eventos con $A=\{X_N=w_N\}$ y
$B=\{X_1=w_1,\dots,X_{N-1}=w_{N-1}\}$, lo cual es lícito porque $P(B)>0$ por
hipótesis:
\[
P(w_1\cdots w_N) = P(w_1\cdots w_{N-1})\cdot P(w_N\mid w_1\cdots w_{N-1}).
\]
Sustituyendo la hipótesis inductiva en el primer factor se obtiene el producto
de $N$ términos. \qedhere
\end{proof}

\begin{observacion}
\textbf{Puntos a verificar contra el material del curso.} Primero, la convención
sobre el producto vacío en el caso base. Segundo, y más importante, qué se hace
cuando algún prefijo tiene probabilidad cero: la convención usual es declarar
que la conjunta es cero, pero eso es una convención y no un teorema, ya que el
factor correspondiente sencillamente no está definido.
\end{observacion}

\subsection{El orden de la factorización es una elección}

La regla de la cadena vale para \emph{cualquier} orden de los eventos: la
demostración anterior funciona igual reetiquetando índices, y la conjunta
resultante es idéntica. Factorizar de izquierda a derecha es una decisión de
modelado, no una necesidad matemática.

Esto tiene una consecuencia que se aprecia mejor en la sección siguiente: la
factorización exacta es simétrica respecto al orden, pero el supuesto de Markov
\emph{no} lo es. Truncar el historial a los $n-1$ símbolos anteriores no es lo
mismo que truncarlo a los $n-1$ posteriores. La aproximación rompe una simetría
que la identidad preservaba.

\section{Cadenas de Markov y el supuesto de Markov}

\subsection{Tres objetos distintos}

Conviene separar tres cosas que suelen mezclarse.

Un \emph{proceso estocástico} es simplemente una secuencia de variables
aleatorias $X_1,X_2,\dots$ sobre un espacio de estados $S$. En un modelo de
lenguaje sobre palabras, $X_i$ es la palabra en la posición $i$ y $S=\Vocab$; en
un etiquetador morfosintáctico, $X_i$ es la etiqueta y $S$ el conjunto de
categorías. Sin supuesto alguno, describir la conjunta de un proceso arbitrario
de longitud $N$ requiere del orden de $|S|^N$ números.

\begin{definicion}
Un proceso es \emph{markoviano de orden $k$} si para todo $i$ y toda asignación
de valores
\[
P(X_i \mid X_1,\dots,X_{i-1}) \;=\; P(X_i \mid X_{i-k},\dots,X_{i-1}).
\]
Cuando se dice ``cadena de Markov'' sin calificar, se sobreentiende $k=1$.
\end{definicion}

\begin{definicion}
Una \emph{cadena de Markov} (de orden $1$) sobre un espacio finito es la
tripleta $(S,\pi,A)$ donde $S$ es el conjunto de estados, $\pi$ es la
distribución inicial con $\sum_{s}\pi(s)=1$, y $A$ es la matriz de transición
con $A_{ss'}=P(X_{i+1}=s'\mid X_i=s)$, sujeta a
\[
\sum_{s'\in S} A_{ss'} = 1 \qquad\text{para cada } s\in S .
\]
La cadena es \emph{homogénea} si $A$ no depende de la posición $i$.
\end{definicion}

Obsérvese con cuidado la condición sobre $A$: lo que suma $1$ es cada
\emph{renglón}, no cada columna. Es exactamente la misma observación sobre el
espacio de normalización que reaparecerá en la sección \ref{sec:mle}, ahora en
forma matricial.

\begin{ejemplo}[Un ejemplo mínimo, para fijar la mecánica]
Estados $S=\{\text{Sol},\text{Lluvia}\}$ con
\[
A=\begin{pmatrix} 0.8 & 0.2 \\ 0.4 & 0.6 \end{pmatrix},
\]
donde el renglón indica el estado de hoy y la columna el de mañana. Cada renglón
suma $1$; las columnas no suman nada en particular ($0.8+0.4=1.2$). Con
$\pi(\text{Sol})=1$,
\[
P(\text{Sol},\text{Sol},\text{Lluvia}) = \pi(\text{Sol})\cdot
A_{\text{Sol},\text{Sol}}\cdot A_{\text{Sol},\text{Lluvia}} = 1\cdot 0.8\cdot
0.2 = 0.16 .
\]
La probabilidad de una secuencia es el producto de las transiciones a lo largo
del camino, multiplicado por la inicial.
\end{ejemplo}

\subsection{El supuesto aplicado al lenguaje}

Un modelo de $n$-gramas resulta de imponer la propiedad de Markov de orden $n-1$
sobre la factorización \eqref{eq:chain}:
\begin{equation}\label{eq:markov}
P(w_i \mid w_1\cdots w_{i-1}) \;\approx\; P(w_i \mid w_{i-n+1}\cdots w_{i-1}).
\end{equation}
Con $n=1$ se tiene el unigrama (historial de cero palabras), con $n=2$ el
bigrama (historial de una), con $n=3$ el trigrama (historial de dos). Nótese que
$n$ cuenta el bloque completo, historial \emph{más} la palabra predicha; el
historial tiene longitud $n-1$.

\begin{idea}[title={Qué se pierde, exactamente}]
El modelo induce una relación de equivalencia sobre el conjunto de historiales:
\[
u \sim u' \iff \text{los últimos } n-1 \text{ símbolos de } u \text{ y de } u'
\text{ coinciden,}
\]
y opera sobre las clases de equivalencia, no sobre los historiales. Toda
distinción interna a una clase es \textbf{inaccesible}: no atenuada ni
aproximada, sino irrecuperable. La contrapartida es que la conjunta pasa de
requerir $|\Vocab|^N$ parámetros a requerir $|\Vocab|^{n}$, cifra independiente
de $N$.
\end{idea}

\subsection{El orden depende de cómo se defina el estado}

Un trigrama es markoviano de orden $2$ sobre palabras, lo cual parece salirse
del marco de las cadenas de orden $1$. No es así, y la construcción que lo
resuelve es estándar.

\begin{proposicion}\label{prop:orden}
Todo proceso markoviano de orden $k$ sobre $S$ es equivalente a una cadena de
Markov de orden $1$ sobre $S^k$, tomando como estado la tupla de los últimos $k$
símbolos.
\end{proposicion}

Para el trigrama el estado es el par $(w_{i-2},w_{i-1})$ y la transición
$(u,v)\to(v,w)$ tiene probabilidad $P(w\mid u,v)$. Obsérvese la restricción
estructural: desde el estado $(u,v)$ sólo son alcanzables estados cuya primera
componente es $v$, de modo que la matriz resultante es enormemente dispersa.

\begin{idea}[title={Consecuencia conceptual}]
El orden \textbf{no es una propiedad del proceso} sino de cómo se definió el
espacio de estados. ``Orden $1$'' no significa ``memoria de un paso'' en ningún
sentido absoluto: una cadena de orden $1$ cuyos estados son tuplas de diez
palabras recuerda diez palabras. La memoria está codificada en el estado. Lo
único que no es gratis es el tamaño: el espacio pasa de $|\Vocab|$ a
$|\Vocab|^{k}$.
\end{idea}

\subsection{El lenguaje no es markoviano de ningún orden finito}

La afirmación precisa es la siguiente: no existe ningún $k$ finito tal que, para
toda oración del español, $w_i$ sea condicionalmente independiente de $w_1\cdots
w_{i-k-1}$ dado $w_{i-k}\cdots w_{i-1}$.

Su forma operativa es un cuantificador universal con contraejemplo constructivo:
\emph{para todo $k$ que se proponga, existe un par de oraciones que comparten
sus últimos $k$ tokens y exigen continuaciones distintas}. El procedimiento de
construcción, en español, puede basarse en la selección de modo verbal:
\begin{center}
\begin{tabular}{ll}
\emph{Espero que mañana} \underline{\hspace{1.6cm}} & $\to$ \emph{llueva}
(subjuntivo)\\
\emph{Sé que mañana} \underline{\hspace{1.6cm}} & $\to$ \emph{lloverá}
(indicativo)
\end{tabular}
\end{center}
Dado cualquier $k$, basta insertar material intermedio hasta empujar el
disparador (\emph{espero} frente a \emph{sé}) fuera de la ventana: \emph{Espero
que mañana por la tarde, si el clima no cambia demasiado y las condiciones se
mantienen como hasta ahora,} \underline{\hspace{1.6cm}}. La concordancia de
número admite construcciones análogas (\emph{El libro que compré para mis
hermanas} frente a \emph{Los libros que compré para mis hermanas}).

Es importante ver por qué el argumento debe tener esta forma. Observar que ``una
sola palabra de contexto no basta'' descarta únicamente $k=1$; por la
proposición \ref{prop:orden}, eso se arreglaría redefiniendo el estado como una
tupla más larga. Lo que hay que descartar es todo $k$ finito simultáneamente.

\begin{observacion}
\textbf{Punto a verificar.} Este argumento está emparentado con la demostración
clásica de que el inglés no es un lenguaje regular (dependencias anidadas
arbitrariamente profundas), pero son afirmaciones distintas: una es sobre
generación entendida como conjunto de cadenas, la otra sobre independencia
condicional en una distribución. Conviene contrastar la formulación exacta con
el material del curso antes de usarla.
\end{observacion}

Tres precisiones sobre lo que la afirmación \emph{no} dice. Primero, no dice que
el modelo sea inútil: un trigrama captura buena parte de la estructura local y
la pregunta relevante no es si el supuesto es verdadero sino cuánto cuesta la
aproximación en cada tarea. Segundo, y este es el punto central, no dice que más
datos lo arreglen: con corpus infinito el estimador converge a la verdadera
$P(w_i\mid \text{últimos } n-1)$, que existe y está bien definida, pero \emph{no
es igual} a $P(w_i\mid \text{todo el historial})$, sino su marginal, y la
marginal ha borrado el disparador. Se trata de error de \textbf{especificación},
sesgo irreducible, y hay que distinguirlo con cuidado del error de
\emph{estimación} por dispersión de la sección \ref{sec:dispersion}, que sí
mejora con datos y es el que justifica el suavizado. Tercero, no dice que las
arquitecturas neuronales escapen: su ventana de contexto también es finita,
aunque el $k$ efectivo sea de órdenes de magnitud mayor y los historiales
compartan parámetros en vez de ser símbolos atómicos.

\section{Estimación por máxima verosimilitud}\label{sec:mle}

\subsection{El criterio y su solución}

En el criterio aparecen cuatro objetos que se escriben parecido y
significan cosas distintas. Conviene fijar los cuatro antes de la
fórmula, porque casi todo lo que sigue es manipulación algebraica una
vez que se sabe qué es cada cual.

\textbf{Los parámetros $\theta_{v\mid u}$.} Un modelo de bigramas no es
más que una tabla de números, uno por cada par (historial $u$, palabra
$v$). A ese número lo llamamos $\theta_{v\mid u}$. Es el mismo objeto
que $P(v\mid u)$, y el cambio de nombre es deliberado: al escribir $P$
se habla de una probabilidad dada, mientras que al escribir $\theta$ se
la trata como una incógnita que hay que determinar a partir del corpus.
Estimar el modelo es exactamente eso, llenar la tabla. El subíndice se
lee ``la palabra $v$ dado el historial $u$'', de modo que cada renglón
de la tabla corresponde a un historial fijo y debe sumar $1$ sobre
todas las palabras posibles.

\textbf{Los conteos $\Count(\cdot)$.} $\Count(x)$ es el número de veces
que la cadena $x$ aparece en el corpus de entrenamiento. No tiene nada
de estadístico: se recorre el corpus una vez y se lleva una cuenta con
lápiz y papel, o con un diccionario. Así, $\Count(uv)$ cuenta las
apariciones del bigrama $uv$, y $\Count(u)$ las del unigrama $u$. La
sección siguiente hace ese conteo a mano sobre un corpus de tres
oraciones. Los conteos son lo único que el corpus aporta: son un
resumen del corpus y, una vez calculados, el corpus ya no se vuelve a
mirar.

\textbf{La verosimilitud $L(\theta)$.} Es la probabilidad que el modelo
asigna al corpus de entrenamiento, leída como función de los parámetros
y no de los datos. Los datos están fijos, ya se observaron; lo que
varía es la tabla, y $L$ mide qué tan bien cada llenado posible de esa
tabla explica lo que se vio. La probabilidad del corpus, por la regla
de la cadena, es el producto de un factor $\theta_{v\mid u}$ por cada
posición del corpus. Al agrupar las posiciones que comparten el mismo
par $(u,v)$, ese factor idéntico se repite $\Count(uv)$ veces y el
producto se vuelve una potencia. De ahí sale el exponente:
\[
L(\theta) = \prod_{u,v} \theta_{v\mid u}^{\,\Count(uv)} .
\]

\textbf{La log-verosimilitud $\ell(\theta)$.} Es simplemente
$\log L(\theta)$, y el logaritmo convierte el producto en suma:
\[
\ell(\theta) = \sum_{u}\sum_{v} \Count(uv)\,\log\theta_{v\mid u} .
\]
Se trabaja con $\ell$ y no con $L$ por dos razones. La primera es
práctica: $L$ es un producto de miles de números menores que $1$ y se
desborda por defecto en cualquier aritmética de punto flotante, cosa
que a una suma de logaritmos no le pasa. La segunda es que el logaritmo
es estrictamente creciente, así que $L$ y $\ell$ alcanzan su máximo en
el mismo punto, y derivar una suma es mucho más cómodo que derivar un
producto.

Con eso, el criterio se enuncia en un renglón: el estimador de máxima
verosimilitud es el valor de $\theta$ que maximiza $\ell$ sin salirse
de la restricción de normalización,
\[
\hat\theta=\argmax_\theta \ell(\theta)
\qquad\text{sujeto a}\qquad
\sum_v \theta_{v\mid u}=1 \ \text{ para cada } u .
\]
El sombrero de $\hat\theta$ distingue el valor estimado a partir de un
corpus concreto del parámetro $\theta$ como incógnita.

\begin{teorema}\label{thm:mle}
El máximo de $\ell$ sobre el símplex se alcanza en las frecuencias relativas:
\[
\hat{P}(w_i \mid w_{i-n+1}\cdots w_{i-1}) \;=\; \frac{\Count(w_{i-n+1}\cdots
w_{i})}{\Count(w_{i-n+1}\cdots w_{i-1})},
\]
y en particular, para el bigrama, $\hat{P}(v\mid u) = \Count(uv)/\Count(u)$.
\end{teorema}

\begin{proof}
\emph{Paso 1 (separabilidad).} La log-verosimilitud se descompone como
$\ell(\theta)=\sum_u \ell_u(\theta)$ con $\ell_u=\sum_v
\Count(uv)\log\theta_{v\mid u}$. Cada término involucra parámetros disjuntos y
su restricción de normalización es independiente de las demás; por lo tanto
maximizar $\ell$ equivale a maximizar cada $\ell_u$ por separado. Esto no es
evidente a priori: se sigue de la estructura de la factorización
\eqref{eq:chain}. Fíjese un $u$ en adelante.

\emph{Paso 2 (condición de primer orden).} Con multiplicador de Lagrange para la
restricción $\sum_v\theta_{v\mid u}=1$,
\[
\mathcal{L}=\sum_v \Count(uv)\log\theta_{v\mid
u}-\lambda\Big(\sum_v\theta_{v\mid u}-1\Big),
\qquad
\frac{\partial\mathcal{L}}{\partial\theta_{v\mid
u}}=\frac{\Count(uv)}{\theta_{v\mid u}}-\lambda=0,
\]
de donde $\theta_{v\mid u}=\Count(uv)/\lambda$.

\emph{Paso 3 (despeje del multiplicador).} Imponiendo la restricción,
\[
\sum_v \theta_{v\mid u} = \frac{1}{\lambda}\sum_v \Count(uv) = 1
\quad\Longrightarrow\quad
\lambda = \sum_v \Count(uv) = \Count(u),
\]
donde la última igualdad usa el lema \ref{lem:conteos}. Sustituyendo,
$\hat\theta_{v\mid u}=\Count(uv)/\Count(u)$.

\emph{Paso 4 (que es máximo).} La segunda derivada es
$\partial^2\ell_u/\partial\theta_{v\mid u}^2=-\Count(uv)/\theta_{v\mid u}^2<0$,
y las cruzadas son nulas: la hessiana es diagonal con entradas negativas, luego
$\ell_u$ es estrictamente cóncava en $\theta>0$ y el punto crítico es máximo
global sobre el símplex. \qedhere
\end{proof}

\begin{lema}\label{lem:conteos}
$\sum_{v\in\Vocab}\Count(uv)=\Count(u)$ para todo $u$ que sea seguido por algún
símbolo en cada una de sus apariciones.
\end{lema}

La justificación es directa: cada aparición de $u$ va seguida por exactamente
una palabra, de modo que las apariciones de $u$ quedan particionadas según cuál
sea esa palabra. La hipótesis restrictiva no es cosmética, y la sección
siguiente exhibe dónde falla.

\begin{idea}[title={La fórmula es consecuencia, no punto de partida}]
La estimación por frecuencias relativas no es una regla \emph{ad hoc} que
resulte razonable; es la solución del problema de optimización planteado arriba.
Reconocerlo importa porque el suavizado, más adelante, será un \textbf{abandono
deliberado} de este criterio: se acepta peor ajuste al entrenamiento a cambio de
no asignar cero a lo no visto.
\end{idea}

\section{Trazado completo sobre un corpus juguete}

Se toma el corpus de tres oraciones
\begin{center}
\texttt{<s> el perro ladra </s>} \qquad
\texttt{<s> el gato duerme </s>} \qquad
\texttt{<s> el perro duerme </s>}
\end{center}
y se estima el bigrama por máxima verosimilitud, llenando las tablas a mano.

\subsection{Conteos}

Los unigramas, que actuarán como denominadores:

\begin{center}
\begin{tabular}{lccccccc}
\toprule
$u$ & \bos & el & perro & gato & ladra & duerme & \eos \\
\midrule
$\Count(u)$ & 3 & 3 & 2 & 1 & 1 & 2 & 3 \\
\bottomrule
\end{tabular}
\end{center}

Los bigramas, enumerados por oración: de la primera, (\bos, el), (el, perro),
(perro, ladra), (ladra, \eos); de la segunda, (\bos, el), (el, gato), (gato,
duerme), (duerme, \eos); de la tercera, (\bos, el), (el, perro), (perro,
duerme), (duerme, \eos). Agrupando:

\begin{center}
\begin{tabular}{lc@{\qquad}lc}
\toprule
$uv$ & $\Count(uv)$ & $uv$ & $\Count(uv)$\\
\midrule
(\bos, el)      & 3 & (perro, duerme) & 1\\
(el, perro)     & 2 & (gato, duerme)  & 1\\
(el, gato)      & 1 & (ladra, \eos)   & 1\\
(perro, ladra)  & 1 & (duerme, \eos)  & 2\\
\bottomrule
\end{tabular}
\end{center}

\subsection{Condicionales y verificación de la normalización}

\begin{center}
\begin{tabular}{lll}
\toprule
Historial $u$ & Distribución sobre $v$ & $\sum_v \hat P(v\mid u)$\\
\midrule
\bos    & $\hat P(\text{el}\mid \bos)=3/3=1$ & $1$\\
el      & $\hat P(\text{perro}\mid\text{el})=2/3$, \quad $\hat
P(\text{gato}\mid\text{el})=1/3$ & $2/3+1/3=1$\\
perro   & $\hat P(\text{ladra}\mid\text{perro})=1/2$, \quad $\hat
P(\text{duerme}\mid\text{perro})=1/2$ & $1$\\
gato    & $\hat P(\text{duerme}\mid\text{gato})=1/1=1$ & $1$\\
ladra   & $\hat P(\eos\mid\text{ladra})=1/1=1$ & $1$\\
duerme  & $\hat P(\eos\mid\text{duerme})=2/2=1$ & $1$\\
\bottomrule
\end{tabular}
\end{center}

Cada renglón suma $1$, como exige la condición sobre la matriz de transición. La
verificación del lema \ref{lem:conteos}, renglón por renglón, da
$\Count(\text{el})=3=2+1$, $\Count(\text{perro})=2=1+1$,
$\Count(\text{duerme})=2=2$, todos correctos; pero para $u=\eos$ se tiene
$\Count(\eos)=3$ mientras que $\sum_v \Count(\eos\,v)=0$. Ahí falla la
identidad.

\begin{idea}[title={Los símbolos de frontera son asimétricos}]
El lema falla exactamente para los tokens que nunca son seguidos por nada, que
tratando cada oración por separado es únicamente \eos. El arreglo estándar
consiste en no tratar \eos\ como historial: es un estado absorbente y se excluye
del dominio de la matriz de transición. Simétricamente, \bos\ nunca aparece como
segundo elemento de un bigrama. Por lo tanto \textbf{el vocabulario de emisión y
el conjunto de historiales no son el mismo conjunto}, y confundirlos al contar
parámetros o al normalizar es una de las causas más frecuentes de reportar una
distribución que no suma $1$.
\end{idea}

Vale la pena precisar para qué sirve cada símbolo. \bos\ existe para dar
historial a las primeras posiciones: sin él, $\hat P(w_1)$ carecería de
condicionante y el primer factor de la cadena sería de naturaleza distinta a los
demás; para un modelo de orden $n$ se anteponen $n-1$ copias. \eos\ existe
porque es el símbolo emitible que permite que el proceso termine, y sin él la
conjunta no queda definida sobre $\Vocab^*$ sino sobre secuencias de longitud
fija.

\subsection{Probabilidad de las oraciones y qué revela la suma}

Multiplicando a lo largo de cada camino:
\begin{align*}
\hat P(\text{el perro ladra})  &= 1\cdot\tfrac{2}{3}\cdot\tfrac{1}{2}\cdot 1 =
\tfrac13,\\
\hat P(\text{el gato duerme})  &= 1\cdot\tfrac{1}{3}\cdot 1\cdot 1 = \tfrac13,\\
\hat P(\text{el perro duerme}) &= 1\cdot\tfrac{2}{3}\cdot\tfrac{1}{2}\cdot 1 =
\tfrac13,
\end{align*}
cuya suma es exactamente $1$.

La lectura ingenua sería que el modelo está bien normalizado y por tanto es
correcto. Es la lectura equivocada. Que las tres oraciones del entrenamiento
agoten toda la masa significa que el modelo \textbf{no asigna probabilidad
positiva a ninguna otra cadena}. La oración \emph{el gato ladra}, perfectamente
gramatical y compuesta sólo de palabras del vocabulario, recibe
\[
\hat P(\text{el gato ladra}) = 1\cdot\tfrac{1}{3}\cdot 0\cdot 1 = 0,
\]
porque $\Count(\text{gato ladra})=0$. El modelo genera exactamente tres cadenas
con probabilidad positiva, las tres vistas. Es un caso extremo pero limpio de lo
que hace la máxima verosimilitud: maximizar la probabilidad del entrenamiento
hasta dejar todo lo demás en cero.

Obsérvese que aquí la normalización global sí se cumple; lo que está mal no es
cuánta masa hay sino \emph{dónde} se colocó.

\subsection{Contraste con el trigrama}

Sobre el mismo corpus, $\hat P(\text{duerme}\mid\text{el perro}) =
\Count(\text{el perro duerme})/\Count(\text{el perro}) = 1/2$, que coincide con
el bigrama $\hat P(\text{duerme}\mid\text{perro})=1/2$. La coincidencia es un
accidente del corpus: \emph{perro} siempre viene precedido de \emph{el}, así que
el contexto adicional no aporta información alguna.

Más ilustrativo es el conteo de historiales. Los pares con conteo positivo son
(\bos, el), (el, perro), (el, gato), (perro, ladra), (perro, duerme) y (gato,
duerme): seis observados frente a $6^2=36$ posibles con seis tipos de palabra.
Es decir, más del ochenta por ciento de la tabla está vacía en un corpus de tres
oraciones.

\section{Dispersión: por qué la estimación se rompe}\label{sec:dispersion}

Un modelo de $n$-gramas tiene del orden de $|\Vocab|^{n-1}(|\Vocab|-1)$
parámetros libres. Con $|\Vocab|=10^4$ y $n=3$ eso asciende a unos $10^{12}$
valores, cifra que ningún corpus cubre. La consecuencia es inmediata: en el
conjunto de prueba aparecen $n$-gramas con conteo cero, la conjunta se anula por
ser un producto y la log-verosimilitud se va a $-\infty$.

\begin{idea}[title={Dos limitaciones distintas que conviene no mezclar}]
El sesgo por el supuesto de Markov (sección 3) es error de \emph{especificación}
y no mejora con más datos. La dispersión es error de \emph{estimación} y sí
mejora con más datos, aunque nunca lo suficiente. El suavizado ataca la segunda,
no la primera.
\end{idea}

\section{Fuga de masa: normalización local frente a global}\label{sec:fuga}

Que cada renglón de la matriz sume $1$ no garantiza que la conjunta sume $1$
sobre $\Vocab^*$. El siguiente par de ejemplos lo muestra sobre el modelo más
pequeño posible.

\begin{ejemplo}[Con fuga]
Vocabulario de emisión $\{a,\eos\}$, historiales $\{\bos, a\}$, con
\[
P(a\mid\bos)=1,\qquad P(a\mid a)=1,\qquad P(\eos\mid a)=0 .
\]
Cada renglón suma $1$: hay un único evento con masa en cada uno. Sin embargo,
toda cadena finita $a^k\eos$ tiene probabilidad $1\cdot 1^{k-1}\cdot 0 = 0$, de
modo que
\[
\sum_{s\in\Vocab^*} P(s) = 0 \neq 1 .
\]
Toda la masa se fue a la cadena infinita $aaa\cdots$, que no es elemento de
$\Vocab^*$.
\end{ejemplo}

\begin{ejemplo}[Sin fuga]
El mismo modelo con $P(a\mid a)=0.9$ y $P(\eos\mid a)=0.1$ asigna
$P(a^k\eos)=0.9^{\,k-1}\cdot 0.1$ para $k\ge1$, y por serie geométrica
\[
\sum_{k\ge1} 0.9^{\,k-1}(0.1) = \frac{0.1}{1-0.9} = 1 .
\]
La diferencia con el caso anterior es que el ciclo tiene probabilidad de escape
estrictamente positiva.
\end{ejemplo}

\begin{observacion}
\textbf{Punto a verificar, y lo señalo como incertidumbre real.} La condición
general que caracteriza cuándo un modelo de $n$-gramas es \emph{estricto}
(tight, es decir, cuando la conjunta suma $1$ sobre $\Vocab^*$) involucra el
radio espectral de la matriz de transición restringida a los estados no
terminales, y no la enuncio aquí porque no la tengo verificada con la precisión
necesaria. Existe además la \emph{intuición}, no verificada, de que un modelo
estimado por máxima verosimilitud a partir de un corpus de oraciones completas
resulta siempre estricto, ya que todo estado alcanzable tiene un camino
observado hacia \eos; suena razonable y hay resultados análogos para gramáticas
probabilísticas, pero no debe darse por cierto sin contrastarlo. El ejemplo con
fuga de arriba fue construido a mano y no proviene de ninguna estimación por
máxima verosimilitud.
\end{observacion}

\section{Suavizado}

El problema, en una línea: la máxima verosimilitud asigna cero a lo no visto, y
el cero es letal porque la conjunta es un producto. La solución siempre tiene la
misma forma, quitar masa a lo visto y dársela a lo no visto; las familias de
métodos se distinguen en cuánta masa se quita y a dónde se manda.

\subsection{Add-$k$ (Laplace)}

\[
\hat{P}_{\text{add-}k}(v\mid u) = \frac{\Count(uv)+k}{\Count(u)+k|\Vocab|} .
\]
El término $k|\Vocab|$ del denominador es exactamente lo que se necesita para
compensar el $k$ sumado a cada uno de los $|\Vocab|$ numeradores; esa es la
condición de normalización, y verificarla consiste en sumar sobre $v$. Con $k=1$
se obtiene el suavizado de Laplace.

Sobre el corpus juguete, con $|\Vocab|=6$ y $k=1$, resulta $\hat
P(\text{perro}\mid\text{el})=(2+1)/(3+6)=1/3$ frente a $2/3$ de la máxima
verosimilitud, es decir la mitad de la masa original; y $\hat
P(\text{ladra}\mid\text{gato})=(0+1)/(1+6)=1/7$, que antes era cero. El defecto
es visible en esos números: cuando $|\Vocab|$ es grande, el término $k|\Vocab|$
domina el denominador y la estimación se aplana. Add-$k$ redistribuye demasiada
masa, lo cual resulta aceptable en clasificación de texto pero es malo para
modelado de lenguaje.

\subsection{Descuento absoluto}

Se resta una cantidad fija $d\in(0,1)$ a cada conteo positivo:
\[
\hat{P}_d(v\mid u) = \frac{\max\{\Count(uv)-d,\;0\}}{\Count(u)} .
\]
Esto ya no suma $1$. La masa liberada en el contexto $u$ es
\[
\gamma(u) = \frac{d\cdot N_{1+}(u)}{\Count(u)},
\qquad
N_{1+}(u) = |\{v : \Count(uv)>0\}| ,
\]
y es precisamente esa $\gamma(u)$ lo que los métodos siguientes reparten entre
lo no visto, cada uno a su manera.

\begin{observacion}
La justificación empírica del descuento fijo proviene de la observación de que
el conteo esperado en un corpus de prueba de un bigrama que apareció $r$ veces
en entrenamiento es aproximadamente $r-0.75$ para valores medianos de $r$, de
donde el valor típico $d=0.75$. El resultado es estándar (Church y Gale), pero
la cifra conviene verificarla en el material del curso.
\end{observacion}

\subsection{Back-off}

Se recurre al orden menor \emph{sólo} cuando el $n$-grama no se observó:
\[
\hat{P}_{\text{bo}}(v\mid u) =
\begin{cases}
\dfrac{\max\{\Count(uv)-d,0\}}{\Count(u)} & \text{si } \Count(uv)>0,\\[10pt]
\alpha(u)\,\hat{P}_{\text{bo}}(v) & \text{si } \Count(uv)=0 .
\end{cases}
\]
El coeficiente $\alpha(u)$ no se elige sino que se calcula: es la constante que
hace que la masa $\gamma(u)$ liberada por el descuento se reparta exactamente
entre los no vistos, en proporción a la distribución de orden menor, es decir
\[
\alpha(u) = \frac{\gamma(u)}{\sum_{v\,:\,\Count(uv)=0} \hat P_{\text{bo}}(v)} .
\]

\subsection{Interpolación}

Se mezclan los órdenes \emph{siempre}, se haya visto o no el $n$-grama:
\[
\hat{P}_{\text{int}}(w_i \mid w_{i-2}w_{i-1}) = \lambda_1 \hat P(w_i\mid
w_{i-2}w_{i-1}) + \lambda_2 \hat P(w_i\mid w_{i-1}) + \lambda_3 \hat P(w_i),
\qquad \sum_j \lambda_j = 1 .
\]
Los $\lambda_j$ se ajustan sobre un conjunto de desarrollo (held-out) y no sobre
el entrenamiento, donde el óptimo sería $\lambda_1=1$, es decir volver a la
máxima verosimilitud. Una variante estándar los hace depender del contexto,
$\lambda_j(u)$, para confiar más en el orden alto cuando $\Count(u)$ es grande.

\begin{center}
\begin{tabular}{lll}
\toprule
 & Back-off & Interpolación \\
\midrule
Cuándo usa el orden menor & sólo si $\Count(uv)=0$ & siempre \\
$\hat P$ de un $n$-grama visto & depende sólo del orden alto & mezcla los
órdenes \\
Normalización & vía $\alpha(u)$ calculada & automática si $\sum_j\lambda_j=1$ \\
\bottomrule
\end{tabular}
\end{center}

\begin{idea}[title={La distinción operativa}]
En interpolación, dos $n$-gramas vistos con el mismo conteo pueden recibir
probabilidades distintas si sus órdenes menores difieren, porque el orden menor
siempre participa. En back-off no ocurre: si el $n$-grama se vio, el orden menor
no interviene en absoluto.
\end{idea}

\subsection{Kneser-Ney: la idea que lo distingue}

El problema que resuelve puede enunciarse con un ejemplo. La palabra
\emph{Francisco} puede ser frecuente como unigrama y sin embargo aparecer casi
siempre dentro de \emph{San Francisco}. Al hacer back-off tras un contexto no
visto, un unigrama basado en frecuencia le asignaría mucha masa, cuando es un
mal candidato en contexto libre.

La corrección consiste en que la distribución de respaldo no mida frecuencia
sino \emph{versatilidad de contextos}:
\[
P_{\text{cont}}(v) = \frac{|\{u : \Count(uv)>0\}|}{|\{(u',v') :
\Count(u'v')>0\}|},
\]
donde el numerador cuenta cuántos contextos distintos preceden a $v$ y el
denominador el número total de tipos de bigrama. Una palabra con frecuencia alta
pero un único contexto queda penalizada. Kneser-Ney combina esta idea con
descuento absoluto e interpolación, y es el mejor suavizado de $n$-gramas
conocido.

\begin{observacion}
\textbf{Punto a verificar.} La fórmula de $P_{\text{cont}}$ se da con confianza
razonable, pero la versión interpolada y modificada de Kneser-Ney, con
descuentos distintos $d_1, d_2, d_{3+}$ según el conteo, tiene detalles de
normalización que conviene revisar directamente en J\&M antes de usarla
formalmente.
\end{observacion}

\subsection{Dónde puede salir mal}

El descuento absoluto con $d>1$ aplicado sin el $\max\{\cdot,0\}$ produce
numeradores negativos para los $n$-gramas de conteo $1$, es decir masa negativa.
Y cualquier esquema en que $\alpha(u)$ o los $\lambda_j$ se fijen a ojo en lugar
de calcularse rompe la normalización sin que resulte evidente en el código. La
verificación es siempre la misma: fijar un $u$, sumar sobre todo $v$ y comprobar
que da $1$.

\section{Perplejidad}

\subsection{Para qué sirve}

Para comparar modelos de lenguaje entre sí sin ejecutarlos en ninguna
tarea. Es la métrica de evaluación estándar: se entrenan dos modelos,
se evalúan sobre el mismo texto de prueba, y el de menor perplejidad
ajusta mejor.

\subsection{Qué mide}

Cuán sorprendido queda el modelo ante texto que no vio. Concretamente,
el número efectivo de opciones que considera en cada posición.

\begin{itemize}[itemsep=0.1em]
  \item Perplejidad $100$: el modelo duda entre unas $100$ palabras en
        cada paso.
  \item Perplejidad $10$: duda entre unas $10$. Mejor modelo.
  \item Modelo uniforme sobre $|\Vocab|$ palabras: perplejidad
        exactamente $|\Vocab|$. Ese es el piso trivial.
\end{itemize}

\subsection{Definición y forma de cálculo}

\begin{definicion}
La \emph{perplejidad} (perplexity) de un texto de prueba $W=w_1\cdots w_N$ bajo
un modelo $P$ es
\[
\mathrm{PP}(W) = P(w_1\cdots w_N)^{-1/N}
= \Bigg(\prod_{i=1}^{N} P(w_i\mid \text{historial})\Bigg)^{-1/N} .
\]
\end{definicion}

En la práctica se calcula en forma logarítmica para evitar desbordamiento por
defecto (underflow):
\[
\mathrm{PP}(W) = \exp\Bigg(-\frac{1}{N}\sum_{i=1}^{N}\ln
P(w_i\mid\text{historial})\Bigg).
\]

\subsection{Relación con la entropía cruzada}

Definiendo la entropía cruzada por token (cross-entropy) como
\[
H(W) = -\frac{1}{N}\sum_{i=1}^{N} \log_2 P(w_i\mid\text{historial}),
\]
se tiene $\mathrm{PP}(W) = 2^{H(W)}$.

\begin{idea}[title={La base debe ser consistente}]
Si $H$ se mide en bits (logaritmo base $2$), la perplejidad es $2^{H}$; si en
nats (base $e$), es $e^{H}$. La perplejidad resultante es la misma en ambos
casos: lo que cambia es el número al que se llama entropía. Mezclar bases
produce cifras que parecen plausibles y son incorrectas, y es de los errores más
frecuentes.
\end{idea}

\subsection{Interpretación y baseline trivial}

La perplejidad es el factor de ramificación promedio ponderado: el número
efectivo de opciones que el modelo considera en cada posición. El caso de
control lo fija: si el modelo es uniforme sobre $|\Vocab|$ palabras, entonces
$P(w_i\mid\cdot)=1/|\Vocab|$ para todo $i$ y
\[
\mathrm{PP} = \big(|\Vocab|^{-N}\big)^{-1/N} = |\Vocab| .
\]
Es decir, \textbf{la perplejidad de un modelo uniforme es el tamaño del
vocabulario}, y ése es el baseline trivial contra el cual debe compararse
cualquier cifra reportada. Menor perplejidad indica mejor modelo.

\subsection{Los detalles donde se rompe}

Respecto a qué cuenta $N$: si \eos\ se incluye en el producto, y debe incluirse
porque forma parte de la probabilidad de la oración, entonces debe incluirse
también en $N$; incluirlo en uno y no en el otro sesga sistemáticamente la
cifra. En cambio \bos\ normalmente no se cuenta en $N$, porque no se predice
sino que condiciona.

Respecto a $\unk$: dos modelos con vocabularios distintos no son comparables por
perplejidad. Un vocabulario pequeño manda muchas palabras a $\unk$, y $\unk$
resulta fácil de predecir precisamente porque acumula masa, de modo que es
posible bajar la perplejidad de forma arbitraria reduciendo el vocabulario. La
comparación sólo tiene sentido con vocabulario idéntico y el mismo tratamiento
de las palabras fuera de vocabulario.

Finalmente, la perplejidad es una métrica intrínseca: mide ajuste probabilístico
al texto de prueba, no utilidad en ninguna tarea concreta. Suele correlacionar
con el desempeño en traducción o reconocimiento de voz, pero no lo implica. Y si
algún $P(w_i\mid\cdot)$ vale cero, la perplejidad es infinita; por eso
perplejidad y suavizado son inseparables, y no puede evaluarse sobre datos
reales un modelo estimado por máxima verosimilitud sin suavizar.

\section{Resumen del esquema}

Un modelo de lenguaje es una distribución sobre $\Vocab^*$. La regla de la
cadena la factoriza en condicionales sin pérdida alguna, y el orden de la
factorización es una elección. El supuesto de Markov trunca el historial a $n-1$
símbolos: es falso para el lenguaje natural, y no lo arreglan más datos porque
se trata de sesgo de especificación, distinto del error de estimación por
dispersión. La máxima verosimilitud estima las condicionales por frecuencias
relativas, resultado que se deriva y no se postula, y asigna cero a todo lo no
observado. El suavizado abandona deliberadamente ese criterio moviendo masa de
lo visto a lo no visto, y las familias de métodos difieren en cuánta masa y
hacia dónde. La perplejidad evalúa el resultado: es $2^{H}$ con $H$ la entropía
cruzada por token, y sólo resulta comparable bajo vocabulario y segmentación
idénticos.

Quedan como puntos explícitamente pendientes de verificación contra el material
del curso: la condición de estrictez del modelo, la constante del descuento
absoluto y los detalles de normalización de Kneser-Ney modificado.

\end{document}

%######################### END OF MODELOS_N_GRAMAS.TEX #########################
%###############################################################################
